\chapter{Limitations}

\section{Performance Claims}

The repository has benchmark tooling, but it does not yet contain a curated set of large, repeated benchmark results.
Current smoke tests validate functionality and rough benchmark execution.
They should not be interpreted as broad performance conclusions.

Future performance claims should specify:

\begin{itemize}
    \item hardware;
    \item compiler and flags;
    \item dataset;
    \item item count;
    \item number of repetitions;
    \item statistical summary;
    \item exact commit.
\end{itemize}

\section{Thread Safety}

heapx does not provide internal synchronization.
Distinct heaps may be used independently by one thread at a time, but concurrent mutation requires external synchronization.
The internal heap id counter is process-local and not atomic.

\section{Linear contains}

\texttt{heapx\_contains()} is a linear pointer-identity search in all current backends.
It is useful in tests and diagnostics, but callers should not treat it as a fast membership query.

\section{API Scope}

The API is intentionally small and heap-shaped.
It does not currently expose:

\begin{itemize}
    \item meld;
    \item increase-key;
    \item custom allocators;
    \item item destructors;
    \item stable ordering among equal priorities;
    \item iteration over heap contents.
\end{itemize}

These omissions are not accidental.
They keep the project focused on the operations needed to compare the current heap backends.

\section{Allocation Policy}

The library uses standard C allocation functions internally.
This is acceptable for the current project, but some production environments require custom allocation hooks.
Adding such hooks would require careful API design because all backends and node pools would need to follow the same allocation policy.

\section{Amortized Bounds and Allocation Failures}

The documented amortized bounds describe normal successful-operation paths.
When temporary allocation fails during consolidation, pointer-heavy backends preserve correctness through fallback behavior.
This may weaken the ideal structural state for that operation.

This is a reasonable engineering choice, but it should remain documented because it separates logical correctness from ideal amortized restructuring.
